% Part: sets-functions-relations
% Chapter: arithmetization
% Section: cuts
%
\documentclass[../../../include/open-logic-section]{subfiles}

\begin{document}

\olfileid{sfr}{arith}{cuts}
\olsection{له $\Rat$ څخه $\Real$ ته}

په اصل کښې د بشپړتيا خاصيت ښيي چې د حقيقي عددونو د کرښې هر ټکی
$\alpha$ دا کرښه په بشپړ ډول پر دوو برخو وېشي: هغه عددونه چې
$\alpha$ ئې تر ټولو لږ پاسنی حد دے، او هغه عددونه چې $\alpha$ ئې
تر ټولو لوے لاندينی حد دے۔ ډېډېکېنډ وړانديز وکړ چې له ناطقو عددونو
څخه د حقيقي عددونو د \emph{جوړولو} دپاره، حقيقي عددونه يوازې د هغو
\emph{پرېکونونو} په توګه وګڼو چې ناطق عددونه وېشي۔ يعنې $\sqrt{2}$
له هغه \emph{پرېکون} سره يو ګڼو چې ناطق عددونه $< \sqrt{2}$ له
ناطق عددونو $> \sqrt{2}$ څخه بېلوي۔

راځئ دا خبره سمه منظمه کړو۔ که ناطق عددونه پر دوو برخو پرې کړو،
نو يوازې د هغې \emph{لاندېنۍ} برخې په کتلو خپل وېش په يوازيني ډول
پېژندلے شو۔ نو په دقيق ډول لاندې تعريف وړاندې کوو:

\begin{defn}[پرېکون]
يو \emph{پرېکون} $\alpha$ د ناطقو عددونو هره ناتشه خاصه پيلنۍ
برخه ده چې تر ټولو لوے غړے نۀ لري۔ يعنې $\alpha$ هله او يوازې هله
پرېکون دے چې:
\begin{enumerate}
	\item \emph{ناتش، خاص}: $\emptyset \neq \alpha \subsetneq \Rat$
	\item \emph{پيلنۍ}: د ټولو $p,q \in \Rat$ دپاره: که $p < q \in \alpha$ وي، نو $p \in \alpha$
	\item \emph{تر ټولو لوے غړے نۀ لري}: د هر $p \in \alpha$ دپاره داسې $q \in \alpha$ شته چې $p < q$ وي
\end{enumerate}
بيا $\Real$ د پرېکونونو سټ دے۔
\end{defn}

نو اوس ويلے شو چې $\sqrt{2} = \Setabs{p \in \Rat}{p^2 < 2\text{
يا }p < 0}$۔ البته بايد وارزوو چې دا په رښتيا يو \emph{پرېکون}
دے، خو دا کار \olref[check]{sec} ته پرېږدو۔

لکه مخکښې، چې ځينې څيزونه مو تعريف کړل، بيا بايد پر هغو بنسټيزې
تابعې او اړيکې تعريف کړو۔ په يوه اسانه تعريف پيل کوو:
\begin{align*}
	\alpha \leq \beta \text{ هله او يوازې هله چې }\alpha \subseteq \beta
\end{align*}
د مرتبې دا تعريف موږ ته اجازه راکوي چې مرکزي پايله \emph{بيان}
کړو: د پرېکونونو سټ د بشپړتيا خاصيت لري۔ په بشپړ ډول د بيان بڼه
داسې ده۔ که $S$ د پرېکونونو يو ناتش سټ وي چې پاسنی حد لري، نو
$S$ تر ټولو لږ پاسنی حد لري۔ په لا تفصيل: يو پرېکون $\lambda$ شته
چې د $S$ پاسنی حد دے، يعنې $(\forall \alpha \in S)\alpha \subseteq
\lambda$، او $\lambda$ تر ټولو لږ داسې پرېکون دے، يعنې
$(\forall \beta \in \Real)((\forall \alpha \in S)\alpha \subseteq \beta \lif \lambda \subseteq \beta)$۔ اوس د پايلې ثبوت دا دے:

\begin{thm}\ollabel{realcompleteness}
د پرېکونونو سټ د بشپړتيا خاصيت لري۔
\end{thm}

\begin{proof}
فرض کړئ $S$ د پرېکونونو هر ناتش سټ دے چې پاسنی حد لري۔ وټاکئ
$\lambda = \bigcup S$۔
%\Setabs{p \in \Rat}{(\exists \alpha \in S)p \in \alpha}$$
لومړے ادعا کوو چې $\lambda$ يو پرېکون دے:
\begin{enumerate}
\item ځکه چې $S$ ناتش دے، لږ تر لږه يو پرېکون په $S$ کښې دے، نو
$\emptyset \neq \lambda$۔ \textbf{د ژباړې د نسخې سمون
(OLARI-003):} انګرېزي سرچينه دا ګام په تېروتنه د پاسني حد له فرض
سره تړي؛ اړين دليل د همدې سټ ناتش والے دے۔ ځکه چې $S$ د پرېکونونو سټ
دے، $\lambda \subseteq \Rat$۔ ځکه چې $S$ پاسنی حد لري، کوم
$p \in \Rat$ د هر پرېکون $\alpha \in S$ څخه بهر دے۔ نو
$p\notin \lambda$، او له دې امله $\lambda \subsetneq \Rat$۔
\item فرض کړئ $p < q \in \lambda$۔ نو داسې $\alpha \in S$ شته
چې $q \in \alpha$۔ ځکه چې $\alpha$ يو پرېکون دے،
$p \in \alpha$۔ نو $p \in \lambda$۔
\item فرض کړئ $p \in \lambda$۔ نو داسې $\alpha \in S$ شته چې
$p \in \alpha$۔ ځکه چې $\alpha$ يو پرېکون دے، داسې $q \in
\alpha$ شته چې $p < q$۔ نو $q \in \lambda$۔
\end{enumerate}
په دې سره ادعا ثابته شوه۔ برسېره پر دې، په څرګند ډول
$(\forall \alpha \in S)\alpha \subseteq \bigcup S = \lambda$؛
يعنې $\lambda$ د $S$ يو پاسنی حد دے۔ اوس فرض کړئ
$\beta \in \mathbb{R}$ هم پاسنی حد دے؛ يعنې
$(\forall \alpha \in S)\alpha \subseteq \beta$۔ د هر
$p \in \Rat$ دپاره، که $p \in \lambda$ وي، نو داسې
$\alpha \in S$ شته چې $p \in \alpha$، او په دې توګه
$p \in \beta$۔ په عمومي کولو $\lambda \subseteq \beta$۔ نو
$\lambda$ د $S$ \emph{تر ټولو لږ} پاسنی حد دے۔
\end{proof}

نو موږ داسې ډېر څيزونه لرو چې د بشپړتيا خاصيت پوره کوي۔ د دې د
ويلو يوه لار دا ده: زموږ په پرېکونونو کښې هېڅ «تشيال» نشته۔
(نو د ناطقو عددونو پر ځاے د حقيقي عددونو نور «پرېکونونه» اخيستل
به هېڅ په زړه پورې نوي څيزونه رانۀ کړي۔)

بيا بايد پر حقيقي عددونو ځينې عمليې تعريف کړو۔ ناطق عددونه په
حقيقي عددونو کښې په دې ټاکنه ورګډوو چې د هر $p \in \Rat$ دپاره
$p_\Real = \Setabs{q \in \Rat}{q < p}$۔ بيا تعريفوو:
\begin{align*}
	\alpha + \beta &= \Setabs{p + q}{p \in \alpha \land q \in \beta}\\
%	\alpha - \beta &\defis \Setabs{p - q}{p \in \alpha \land q \in \Rat \setminus \beta}\\
	\alpha \times \beta &=
	\Setabs{p \times q}{0 \leq p \in \alpha \land 0 \leq q \in \beta} \cup 0^\mathbb{R} & \text{که }\alpha, \beta \geq 0_\Real
%	\alpha \div \beta &\defis
%	\Setabs{\nicefrac{p}{q}}{p \in \alpha \land q \in \Rat \setminus \beta}  & \text{که }\alpha \geq 0_\Real, \beta > 0_\Real
\end{align*}
\textbf{د ژباړې د نسخې سمون (OLARI-004):} سرچينې ناطق عددونه
حقيقي عددونو ته په لاندې‌ليک نښه ورګډ کړي، خو په دې اتحاد کښې د
حقيقي صفر نښه په پورته‌ليک ليکي۔ د فورمول د برابرۍ دپاره اصلي نښه
ساتل شوې؛ دلته ئې ټاکل شوې معنا د حقيقي صفر پرېکون دے۔ د ضرب د
نورو حالتونو د سمبالولو دپاره لومړے وټاکئ: %چې $0_\Real \times \alpha = 0_\Real = \alpha \times 0_\Real$، او ورزيات کړئ:
\begin{align*}
	-\alpha &= \Setabs{p - q}{p < 0 \land q \notin \alpha}
\end{align*}
او بيا وټاکئ:
\begin{align*}
	\alpha \times \beta &\defis
	\begin{cases}
		\mathord{-}\alpha \times \mathord{-}\beta &\text{که }\alpha < 0_\Real\text{ او }\beta < 0_\Real\\
		\mathord{-}(\mathord{-}\alpha \times \beta) &\text{که }\alpha < 0_\Real \text{ او }\beta > 0_\Real\\
		\mathord{-}(\alpha \times \mathord{-}\beta) &\text{که }\alpha > 0_\Real \text{ او }\beta < 0_\Real
	\end{cases}
\end{align*}
\textbf{د ژباړې د نسخې سمون (OLARI-005):} فعاله انګرېزي سرچينه
هغه د ضرب حالتونه نۀ تعريفوي چې يو عامل صفر او بل منفي وي؛ د صفر
قاعده يوازې په غېر فعاله تبصره کښې پاتې ده۔ نو تر هغې قاعدې پرته
پورته ټوټه‌ييز تعريف بشپړ ضرب نۀ تعريفوي۔ بيا بايد وارزوو چې دا هر
تعريف تل يو پرېکون ورکوي۔ په پاے کښې بايد په يوه اسانه (خو اوږده)
څرګندونه کښې وښيو چې په دې ډول تعريف شوي پرېکونونه کټ مټ لکه څنګه
چې بايد، هماغسې چلېږي۔ خو دا ټول \olref[check]{sec} ته پرېږدو۔
\end{document}
